\documentclass[10pt,a4paper]{amsart}
\usepackage[margin=19mm]{geometry}
\usepackage{amsmath,amssymb,mathtools}
\usepackage[scr=rsfs,cal=euler]{mathalfa}
\usepackage{xfrac}
\usepackage[unicode,hidelinks,bookmarks=false]{hyperref}
\newcommand{\ie}{i.e.\ }
\newcommand{\cf}{cf.\ }
\newcommand{\resp}{respectively\ }
\newcommand{\Subsection}{\subsection}
\providecommand{\nobreakdash}{\nobreak\hbox{-}}
\setlength{\emergencystretch}{2em}
\multlinegap0pt
\usepackage{amssymb}
\usepackage[normalem]{ulem}
\usepackage{enumerate}
\usepackage{bm}
\newtheorem{lem}{Lemma}
\newtheorem{rmk}{Remark}
\newtheorem{thm}{Theorem}
\newcommand{\F}{\mathbb F}
\title{Permutation polynomials with a few terms \\ over finite fields}
\author{Kirpa Garg, Sartaj Ul Hasan, Chunlei Li, Hridesh Kumar, Mohit Pal}
\date{Source: arXiv:2503.20982v3 (CC BY 4.0)}
\begin{document}
\maketitle

\noindent\emph{Source and licence.} Permutation polynomials with a few terms \\ over finite fields. Version arXiv:2503.20982v3 (CC BY 4.0), distributed by arXiv under the Creative Commons Attribution 4.0 International licence, http://creativecommons.org/licenses/by/4.0/, as recorded in the arXiv licence metadata for this exact version and shown on the abstract page https://arxiv.org/abs/2503.20982v3. Full licence notice: \texttt{licenses/arxiv-2503.20982-CC-BY-4.0.txt}.\par\smallskip

\noindent\textit{Editorial scope.} Counted unit: the article's thm T05 (source0.tex lines 297--304). The article's complete original proof of it is delivered with it (source0.tex lines 305--327, one \texttt{proof} environment of 23 lines). No mathematics is written, repaired or re-derived: the statement and the proof are the article's own lines, in the article's own notation, and no retained line is substituted or re-flowed.  Retained uncounted as the source-local context the proof needs: lines 210--217; lines 254--268; lines 270--273.  Nothing was omitted from any retained span, so the package carries no omission record.  No retained line contains a citation, so the wrapper carries no bibliography and no bibliography entry.  No retained line contains a figure environment, an image-inclusion command or drawing code, so no image is delivered and the package declares no asset dependency.  Editorial formatting only, in addition: an AMS \texttt{amsart} preamble that restates the article's own declarations and macros used by the retained lines, namely the environments \texttt{lem}, \texttt{rmk}, \texttt{thm} on separate counters, together with the standard packages those lines need.  The retained environments print in this document as Lemma 1, Remark 1, Theorem 1 in order of appearance, whereas the article prints the same items with the numbering of its own full document.  The complete native source of the article as distributed in the arXiv e-print for this exact version is bundled unchanged as \texttt{sources/175455/source0.tex}.
\begin{lem}\label{L23}
 Let $h(X) \in \F_q[X]$ and $r,s$ are positive integers such that $s \mid (q-1)$. Then the function $f(X)=X^rh(X^s)$ permutes $\F_q$ if and only if both conditions
 \begin{enumerate}
  \item $\gcd(r,s)=1$ and
  \item $X^rh(X)^s$ permutes $\mu_{\frac{q-1}{s}}$,
 \end{enumerate}
are satisfied, where $\mu_{\frac{q-1}{s}}$ is the set of $\frac{q-1}{s}$-th roots of unity in the algebraic closure $\overline \F_q$ of $\F_q$.
\end{lem}

 \begin{equation}\label{C05}
  \begin{cases}
   N_0 &= -{\tilde \beta}\beta^3 (\delta^{3q}-{\tilde \delta}^q),\\
   N_1 &= 3{\tilde \beta}\beta^2(\delta^{2q+1}-{\tilde \delta}^q),\\
   N_2 &= -3{\tilde \beta}\beta (\delta^{q+2}-{\tilde \delta}^q),\\
   N_3 &= {\tilde \beta}(\delta^3-{\tilde \delta}^q),
  \end{cases}
  ~~~\mbox{and}~~~
  \begin{cases}
   D_0 &= -\beta^3 (\delta^{3q}-{\tilde \delta}),\\
   D_1 &= 3\beta^2(\delta^{2q+1}-{\tilde \delta}),\\
   D_2 &= -3\beta (\delta^{q+2}-{\tilde \delta}),\\
   D_3 &= (\delta^3-{\tilde \delta}).
  \end{cases}
 \end{equation}

 \begin{rmk}\label{R05}
  Since $ {\nu} \circ X^3 \circ \rho: \mu_{q+1} \rightarrow \mu_{q+1}$ is a permutation of $\mu_{q+1}$, the equation $$h(X):= D_3X^3+D_2X^2+D_1X+D_0 = 0$$
  has no solution $X \in \mu_{q+1}$.
 \end{rmk}

\begin{thm}\label{T05}
Let $q\equiv 2 \pmod 3$ and $h(X)=D_3X^3+D_2X^2 +D_1X+D_0$, where $D_0, D_1, D_2,$ $D_3$ are as given in Equation~\eqref{C05},
and $\beta, {\tilde \beta}$ be elements in $\mu_{q+1}$ satisfying $1+{\tilde \beta}\beta^3=0$. Moreover, $\delta$ and $\tilde \delta$ in the expressions of $D_i$'s, $0 \leq i \leq 3$ are such that $\delta, \tilde \delta \in \F_{q^2}\setminus\F_{q}$ and $\tilde \delta \not \in \{ \delta^3, \delta^{q+2}, \delta^{2q+1}, \delta^{3q}\}$. Then 
$$
f(X)=X^3h(X^{q-1}) = D_3X^{3q}+D_2X^{2q+1} +D_1X^{q+2}+D_0X^3
$$ is a PP over $\F_{q^2}$.

\end{thm}

\begin{proof}
From Lemma~\ref{L23}, we know that $f$ is a permutation of $\F_{q^2}$ if and only if $\gcd(3,q-1)=1$ and $g(X):=X^3h(X)^{q-1}$ permutes $\mu_{q+1}$. Since $q\equiv 2 \pmod 3$, the first condition holds trivially. Therefore it is sufficient to show that if $1+{\tilde \beta}\beta^3=0$ then $g$ permutes $\mu_{q+1}$. From Remark~\ref{R05}, we know that $h(X)=0$ has no solution in $\mu_{q+1}$. Now for any $\alpha \in \mu_{q+1}$, consider
 \[
  g(\alpha)= \alpha^3h(\alpha)^{q-1}=\alpha^3\frac{h(\alpha)^q}{h(\alpha)}= \frac{D_0^q\alpha^3+D_1^q\alpha^2 +D_2^q\alpha+D_3^q}{D_3\alpha^3+D_2\alpha^2 +D_1\alpha+D_0}.
 \]
Let $\displaystyle {\tilde g}(X):= \frac{D_0^qX^3+D_1^qX^2 +D_2^qX+D_3^q}{D_3X^3+D_2X^2 +D_1X+D_0} \in \F_{q^2}(X)$. It is easy to observe that $g$ permutes $\mu_{q+1}$ if and only if the rational function ${\tilde g}(X)$ permutes $\mu_{q+1}$. Consider the following system of equations
\begin{equation} \label{SEE1}
  \begin{cases}
   N_3 &= D_0^q,\\
   N_2 &= D_1^q,\\
   N_1 &= D_2^q,\\
   N_0 &= D_3^q,
  \end{cases}
  \iff
  \begin{cases}
   (\delta^3-{\tilde \delta}^q)({\tilde \beta} +\beta^{3q})=0 ,\\
   (\delta^{q+2}-{\tilde \delta}^q)(3{\tilde \beta}\beta +3\beta^{2q})=0,\\
   (\delta^{2q+1}-{\tilde \delta}^q)(3{\tilde \beta}\beta^2+3\beta^q)=0,\\
   (\delta^{3q}-{\tilde \delta}^q)(1+{\tilde \beta}\beta^3)  = 0,
  \end{cases}
 \end{equation}
 where $N_0, N_1, N_2$ and $N_3$ are as given in Equation~\eqref{C05}. If all the four conditions in the system~\eqref{SEE1} are satisfied then ${\tilde g}={\nu} \circ X^3 \circ \rho$. It is easy to verify that if $1+{\tilde \beta}\beta^3=0$ then  all the four conditions of the system~\eqref{SEE1} are satisfied. Thus, if $1+{\tilde \beta}\beta^3=0$ then ${\tilde g}$ is a permutation of $\mu_{q+1}$ and consequently $f$ is a permutation of $\F_{q^2}$.
\end{proof}

\end{document}
